\documentclass{article}
\usepackage{helvet}

% \usepackage{fouriernc}
\usepackage{fourier}
\usepackage{tgschola}

%%%%%%%%%%%%%%%%%%%%%%%%%%%%%%%%%%%%%%%%%%%%%%%%%%%%%%%%%%%%%%%%%%%%%%%%%%%%%%%%
%% Thesis meta-information
%%

%% The title of the thesis:
\title{\bf Deep learning for grouped data}

%% The full name of the author (e.g.: James Smith):
\author{Dan Andrei Iliescu}

%% Declaration date:
\date{3 July 2024}


\begin{document}

\maketitle

\begin{abstract}
    This dissertation explores the problems inherent in applying deep learning algorithms to groups of data. My claim is that groups should be represented as random variables whose values should be inferred from data. This approach has the potential to unlock solutions in many important domains of machine learning, including disentangling the generative factors of data, performing missing data imputation, or training robust predictors. However, grouped data also comes with challenges, especially when the data is high-dimensional and non-linear. Addressing these limitations is the focus of the technical contributions of my doctorate.
\end{abstract}

\end{document}
